\section{Restoring bit circuits with a raw normalized readout}
\label{sec:raw-restoring-bits}
This comparison refinement separates an analog prefix from later digital
computation. It concerns the same gamma-gated Student hidden prior and
normalized Gaussian readout. No additional restriction is imposed on
the learner's prior.

Fix an input law $P_X$ and a measurable good set $\mathcal A$ with
$P_X(\mathcal A^c)\le\xi$. A teacher prefix has depth $J$, width at most
$m$, coefficient cap $B_0$, and resource $Q_{\rm pre}$ in the neuron-gate
comparison. On $\mathcal A$, the final features used by the digital block approximate specified
bits in $\{-1,1\}$ within $\delta/4$. Other final prefix features may be retained but have zero intended outgoing coefficients in that block.
After the prefix, let a layered signed-bit circuit use unweighted
majority gates, including copying and negation. Before multiplying by
a gain $k$, each gate has integer coefficients of magnitude at most
one, a bias of magnitude at most $m+1$, and signed preactivation margin
at least one on its intended input bits. Let $Q_{\rm dig}$ count all
active dense coefficient positions, hidden vertices and layers after
the prefix, including zero positions between selected neurons.
The actual tanh implementation multiplies each gate's coefficients
and bias by
\[
 k\ge\log(8/\delta),\qquad
 0<\delta\le\frac1{16(m+1)}.
\]
Let the ideal bit readout have coefficients $v_j$, bias $b$, and
$L=\sum_j|v_j|$, with $L+|b|\le4$. Let $g_*$ denote that ideal readout
on $\mathcal A$, extended arbitrarily with $|g_*|\le4$ outside it.

\begin{lemma}[A depth-stable raw bit comparison]
\label{lem:raw-bit-comparison}
For every prescribed finite width $N\ge m$, there is a
sample-independent comparison law $\rho$ for the raw-output learner
such that
\[
 \rho\{\|z_w\|_\infty\le6\}=1,\qquad
 \E_\rho\|z_w-g_*\|_{L^2(P_X)}^2
       \le C(\delta^2+\xi),
\]
and
\begin{align}
 \KL(\rho\Vert\mu_N)\le C Q_{\rm pre}
 &\left[(J+2)\log\{C(d+m+2)(B_0+2)\}
                   +\log(1/\delta)\right]\notag\\
 &+C(Q_{\rm dig}+m+1)
       \log\frac{C(Q_{\rm pre}+Q_{\rm dig}+m+2)(k+1)}{\delta}
       +C(L^2+b^2).
 \label{eq:raw-bit-resource}
\end{align}
Constants depend on the fixed prior parameters but not on learner
width, prefix depth, or digital depth. In particular the digital
block pays its actual parameter count times logarithms, without
a further factor of its depth.
\end{lemma}
\begin{proof}
First construct a comparison before imposing three good events.
Choose the usual disjoint per-vertex location groups at every layer.
In the prefix use the standard gate comparison with scalar precision
small enough that the maximum expected squared uniform error over
all its final selected features is at most
$\delta^2/[10^4(m+1)]$. Its depth-uniform stability bound shows that
the logarithm of this precision is at most
\[
 C(J+2)\log\{C(d+m+2)(B_0+2)\}
       +C\log((m+1)/\delta).
\]
The first good event makes every such feature error at most
$\delta/4$. A union bound and Markov's inequality give this event
probability at least $0.99$. On $\mathcal A$ the actual prefix output
is then within $\delta/2$ of its intended bits. Outgoing gates of
the last prefix layer can be assigned at the next step, because they
do not enter their own features.

At every digital layer put selected gates uniformly near one,
condition the total inactive gate mass to be at most $\epsilon$,
and put selected edge/bias factors uniformly within $\epsilon$ of
their intended gain-$k$ values. All actual tanh features are bounded
by one. The selected coefficient perturbation in one row is
deterministically at most $C m(k+1)\epsilon$. Its inactive background
has expected uniform absolute value at most $C\epsilon$, since the
total inactive gate mass is at most $\epsilon$ and the native factor
has a fixed first absolute moment. These bounds do not require the
preceding features to be independent of one another.

Let $V_{\rm dig}$ be the number of digital gates and choose
\[
 \epsilon\le
 \frac{c}{(V_{\rm dig}+1)(m+1)(k+1)}.
\]
The sum of the expected uniform absolute perturbations over every
digital row is then at most $k/1000$. The second good event makes
each such perturbation at most $k/8$; Markov's inequality gives
failure probability at most $0.01$.

On the first two good events, the digital output error remains
at most $\delta$ at every depth. Indeed, if incoming features are
within $\delta$ of the intended bits, their contribution to the
preactivation error is at most $km\delta\le k/16$.
The ideal signed margin is at least $k$, and the new row perturbation
is at most $k/8$. Thus the actual preactivation has the correct sign
and magnitude at least $3k/4$. Its tanh output differs from that sign
by at most $2e^{-3k/2}\le\delta$. Induction proves the assertion,
starting from the prefix error $\delta/2$.

At the normalized Gaussian readout use the variation-adapted
Dirichlet comparison: selected latent gate centers are
\[
 \omega_j=(|v_j|+L/r)/(2L)
\]
when $L>0$, where $r$ is the number of selected output bits; use
$\omega_j=1/r$ if $L=0$. They sum to one and are at least $1/(2r)$.
The selected Gaussian centers are $v_j/\sqrt{\omega_j}$.
Their quadratic cost is at most $2L^2$, and the bias adds $Cb^2$.
Choose its scalar precision at most $c\delta/(m+1)$.
The selected effective coefficient errors then have total absolute
size at most $C\delta$, while the inactive contribution has expected
uniform absolute size at most $c\delta$ by the bound
$\sum_{\rm inactive}\sqrt{P_j}\le C\epsilon$ in the readout comparison.
The third good event bounds that inactive absolute contribution by
$\delta/100$; choosing its precision constant small makes its
failure probability at most $0.01$.

The intersection $E$ of the three good events has probability at
least $0.97$ under the comparison law. On $E$, the raw output is
uniformly bounded by six, and on $\mathcal A$ its discrepancy from
$g_*$ is at most $C\delta$. On $\mathcal A^c$, both functions have
absolute value at most six. This gives the stated integrated error
after conditioning the comparison on $E$.

Before conditioning, every selected latent gate has the usual
$\log N$ entropy cost. The disjoint location mixture cancels all
these costs, including those for the final Dirichlet gates.
The prefix pays the first term of
\eqref{eq:raw-bit-resource}. Digital factors, gates and biases pay
$C Q_{\rm dig}\log[C(Q_{\rm pre}+Q_{\rm dig}+m+2)(k+1)/\delta]$;
the selected final factors and gates pay at most the same logarithm
times $C(m+1)$, plus $C(L^2+b^2)$.
The selected marker intervals remain disjoint from inactive gates
because the precisions are chosen smaller if necessary.

Finally conditioning a law $\nu$ on an event of probability at
least $0.97$ increases its KL against the native prior by at most
a fixed factor plus a numerical constant. Explicitly the KL chain
rule for $(E,E^c)$ implies
\[
 \nu(E)\KL(\nu(\cdot\mid E)\Vert\mu_N)
 \le\KL(\nu\Vert\mu_N)+h(\nu(E)),
\]
where $h$ is binary entropy. This proves
\eqref{eq:raw-bit-resource}. Conditioning occurs only in the
comparison proof, not in the actual learner.
\end{proof}

\begin{corollary}[A growing readout range]
\label{cor:raw-bit-range}
If the ideal bit readout instead has $L+|b|\le4M$ for a deterministic
$M\ge1$, the same construction gives
\[
 \rho\{\|z_w\|_\infty\le6M\}=1,\qquad
 \E_\rho\|z_w-g_*\|_{L^2(P_X)}^2
       \le C M^2(\delta^2+\xi).
\]
In the resource \eqref{eq:raw-bit-resource}, add $\log(M+1)$ inside
the logarithmic precision terms and retain the exact final cost
$C(L^2+b^2)\le C M^2$.
\end{corollary}
\begin{proof}
The digital and prefix comparison is unchanged. At the readout use
precision $c\delta/[(m+1)(M+1)]$, the selected centers
$v_j/\sqrt{\omega_j}$ for the enlarged coefficients, and the same
constant-probability event bounding its inactive absolute contribution
by $\delta/100$. Selected coefficient errors have total absolute
size $C\delta$, while the bit approximation contributes at most
$L\delta\le4M\delta$. The selected absolute-sum norm plus bias is
at most $4M+C\delta$, and the inactive term is at most $\delta/100$,
giving the support bound after reducing the precision constant.
The complement of the good input set contributes at most $CM^2\xi$
to integrated squared error. Only the precision logarithm and the
Gaussian selected-center cost change.
\end{proof}

\paragraph{Interface to the raw-tail theorem.}
The raw-tail statistical proof only needs the comparison supported
inside its auxiliary event
$E_n=\{\|z_w\|_\infty\le M_n\}$ and a population-risk comparison.
For all sufficiently large $n$, $M_n\ge6$, so the law above is already
supported there. Moreover
\[
 \KL(\rho\Vert\mu_N(\cdot\mid E_n))
   =\KL(\rho\Vert\mu_N)+\log\mu_N(E_n)
   \le\KL(\rho\Vert\mu_N).
\]
If $\delta\le n^{-2}$ and $\xi\le n^{-4}$, Cauchy--Schwarz and
$\E Y^2<\infty$ give
$\E_\rho R(z_w)\le R(g_*)+C(1+\sqrt{\E Y^2})n^{-2}$.
The remaining raw-tail comparison between the auxiliary and actual
posterior is unchanged. Thus the resource above can replace the
generic depth-amplified comparison resource in the raw oracle.
For a quantizer prefix of logarithmic depth and logarithmic-bit
output, extra bit-copying layers cost only their selected dense
block count times logarithms, typically $H\,\operatorname{polylog}n$.

For a growing range $M$, the same interface applies once the auxiliary
raw-tail cutoff satisfies $M_n\ge6M$. Choosing
$\delta\le n^{-3}/M^2$ and $\xi\le n^{-6}/M^4$ gives an
$O((1+\sqrt{\E Y^2})n^{-2})$ comparison-risk remainder, since
$|g_*|\le4M$. In particular $M=O(\sqrt{\log n})$ only changes the
compiler resources and statistical remainder by logarithmic factors.
The population comparison itself only needs a finite second moment;
the imported raw-tail oracle still requires its stated marginal
sub-Gaussian label assumption.

\begin{corollary}[The power-readout family]
\label{cor:raw-bit-power}
Fix $p>1$. Replace the final readout by
$G_0+\sum_jP_j^{1-1/p}G_jh_j^H$, with independent factor and bias
density proportional to $e^{-|G|^p}$.
The preceding comparison conclusions remain true, with constants
depending on $p$, and the final magnitude cost becomes
\[
 C_p\left[\left(\sum_j|v_j|\right)^p+|b|^p\right]
 \le C_p M^p.
\]
All other resource terms change only by constants and logarithmic
precision factors.
\end{corollary}
\begin{proof}
Put $\alpha=1-1/p$. Use the same selected gate centers $\omega_j$,
but the selected factor center is $v_j/\omega_j^\alpha$.
For $L>0$, $\omega_j\ge|v_j|/(2L)$, so
\[
 \sum_j\frac{|v_j|^p}{\omega_j^{p-1}}
 \le(2L)^{p-1}\sum_j|v_j|
 =2^{p-1}L^p.
\]
For $L=0$ all selected centers are zero. The power-readout comparison
conditions the inactive latent gate sum to be at most
$\epsilon^{1/\alpha}$ and then bounds
$\sum_{\rm inactive}S_j^\alpha\le C_{\vartheta,p}\epsilon$
at constant additional entropy. The identity
\[
 \E\left[\sum_{\rm inactive}S_j^\alpha
          \,\middle|\,\sum_{\rm inactive}S_j=t\right]
 =
 t^\alpha(N-r)
 \frac{\Gamma(a+\alpha)\Gamma(b_0)}
      {\Gamma(a)\Gamma(b_0+\alpha)}
 \le C_{\vartheta,p}t^\alpha
\]
proves that assertion, with $a=\vartheta/N$ and
$b_0=(N-r)a$; the empty inactive set is interpreted separately.
The total gate sum remains close to one, so selected effective
coefficient errors have total size at most $C_p\delta$ at precision
$c_p\delta/[(m+1)(M+1)]$. Native factors have finite first moments,
so the same inactive-output good event applies. All digital steps
are unchanged. The designated-factor entropy is the displayed
$p$-power cost plus logarithms, and the same constant-probability
conditioning finishes the proof.
\end{proof}
